\section{符号说明}
\begin{table}[H]
\centering
\begin{tabularx}{\textwidth}{CL}
\toprule
符号 & 含义 \\
\midrule
$N$ & 板凳总节数，$N=223$ \\
$t$ & 时间，单位为s \\
$p$ & 螺距，单位为m \\
$v_H$ & 龙头前把手速度，单位为m/s \\
$P_i(t)$ & 第$i$个把手在时刻$t$的平面坐标 \\
$\ell_H$ & 龙头板前后把手中心距，$2.86$ m \\
$\ell_B$ & 普通板和龙尾板前后把手中心距，$1.65$ m \\
$r(\theta)$ & 等距螺线的极坐标半径函数 \\
$S(\theta)$ & 螺线从起点到参数$\theta$的弧长 \\
$R_T$ & 调头空间半径，$4.5$ m \\
$T_{\mathrm{stop}}$ & 问题二的无碰撞终止时刻 \\
$L_{\mathrm{turn}}$ & 调头曲线长度 \\
$v_{\max}$ & 问题五的龙头最大恒定速度 \\
\bottomrule
\end{tabularx}
\end{table}

其中，$P_i(t)$既可表示把手中心，也可在需要时表示与其绑定的板凳参考点，具体含义在各问题中保持一致。本文所有长度单位均为m，速度单位均为m/s，角度单位均为rad。
